% GENERATED FILE. Edit canonical lessons, then run npm run generate:books.
% canonical-source-hash: fnv1a64:3ca84844ed0a8a15
% canonical-lessons: SA-C238-paricarika, SA-C238-abhiyanta, SA-C238-vaijnanikah, SA-C238-araksakah, SA-C238-calakah

\chapter{Nurse, Engineer, Scientist, Policeman, Driver}
\label{ch:sa-nurse-engineer-scientist-policeman-driver}
\hlchaptermodality{238}

\begin{chapteropening}
\textbf{By the end of this chapter:} \emph{I can say a nurse, an engineer, a scientist, a policeman and a driver.}

Complete the last lesson of chapter 238: I can say a nurse, an engineer, a scientist, a policeman and a driver.
\end{chapteropening}

\section[\texorpdfstring{\sk{परिचारिका} (paricārikā)}{परिचारिका (paricārikā)}]{\sk{परिचारिका} (paricārikā) — a nurse}

\label{lesson:SA-C238-paricarika}



\begin{quote}
Before the new one: say the Sanskrit for a weaver, then the Sanskrit for a washerman.
\end{quote}

\subsection*{You'll want to know: \sk{परिचारिका}}
\textbf{\sk{परिचारिका}} — \emph{paricārikā} — \textquotedblleft{}a nurse\textquotedblright{}.

\begin{grammarlens}[title={What you've built}]
One more word for this chapter.
\end{grammarlens}

\subsection*{Your turn}
\begin{itemize}
\raggedright
  \item \emph{Say it:} \emph{paricārikā}
  \item \emph{Say it:} \emph{paricārikā}, once more
  \item \emph{Say it:} say \emph{paricārikā}
  \item \emph{Recall:} say the Sanskrit for a weaver, then the Sanskrit for a washerman, then say \emph{paricārikā} again
\end{itemize}

\begin{tcolorbox}[breakable,colback=teal!4,colframe=teal!35!black,title={Before you move on}]
What does \sk{परिचारिका} mean? (\textquotedblleft{}A nurse\textquotedblright{}.) Say it once more.
\end{tcolorbox}
\section[\texorpdfstring{\sk{अभियन्ता} (abhiyantā)}{अभियन्ता (abhiyantā)}]{\sk{अभियन्ता} (abhiyantā) — an engineer}

\label{lesson:SA-C238-abhiyanta}



\begin{quote}
Before the new one: say the Sanskrit for a washerman, then the Sanskrit for a nurse.
\end{quote}

\subsection*{You'll want to know: \sk{अभियन्ता}}
\textbf{\sk{अभियन्ता}} — \emph{abhiyantā} — \textquotedblleft{}an engineer\textquotedblright{}.

\begin{grammarlens}[title={What you've built}]
One more word for this chapter.
\end{grammarlens}

\subsection*{Your turn}
\begin{itemize}
\raggedright
  \item \emph{Say it:} \emph{abhiyantā}
  \item \emph{Say it:} \emph{abhiyantā}, once more
  \item \emph{Say it:} say \emph{abhiyantā}
  \item \emph{Recall:} say the Sanskrit for a washerman, then the Sanskrit for a nurse, then say \emph{abhiyantā} again
\end{itemize}

\begin{tcolorbox}[breakable,colback=teal!4,colframe=teal!35!black,title={Before you move on}]
What does \sk{अभियन्ता} mean? (\textquotedblleft{}An engineer\textquotedblright{}.) Say it once more.
\end{tcolorbox}
\section[\texorpdfstring{\sk{वैज्ञानिकः} (vaijñānikaḥ)}{वैज्ञानिकः (vaijñānikaḥ)}]{\sk{वैज्ञानिकः} (vaijñānikaḥ) — a scientist}

\label{lesson:SA-C238-vaijnanikah}



\begin{quote}
Before the new one: say the Sanskrit for a nurse, then the Sanskrit for an engineer.
\end{quote}

\subsection*{You'll want to know: \sk{वैज्ञानिकः}}
\textbf{\sk{वैज्ञानिकः}} — \emph{vaijñānikaḥ} — \textquotedblleft{}a scientist\textquotedblright{}.

\begin{grammarlens}[title={What you've built}]
One more word for this chapter.
\end{grammarlens}

\subsection*{Your turn}
\begin{itemize}
\raggedright
  \item \emph{Say it:} \emph{vaijñānikaḥ}
  \item \emph{Say it:} \emph{vaijñānikaḥ}, once more
  \item \emph{Say it:} say \emph{vaijñānikaḥ}
  \item \emph{Recall:} say the Sanskrit for a nurse, then the Sanskrit for an engineer, then say \emph{vaijñānikaḥ} again
\end{itemize}

\begin{tcolorbox}[breakable,colback=teal!4,colframe=teal!35!black,title={Before you move on}]
What does \sk{वैज्ञानिकः} mean? (\textquotedblleft{}A scientist\textquotedblright{}.) Say it once more.
\end{tcolorbox}
\section[\texorpdfstring{\sk{आरक्षकः} (ārakṣakaḥ)}{आरक्षकः (ārakṣakaḥ)}]{\sk{आरक्षकः} (ārakṣakaḥ) — a policeman}

\label{lesson:SA-C238-araksakah}



\begin{quote}
Before the new one: say the Sanskrit for an engineer, then the Sanskrit for a scientist.
\end{quote}

\subsection*{You'll want to know: \sk{आरक्षकः}}
\textbf{\sk{आरक्षकः}} — \emph{ārakṣakaḥ} — \textquotedblleft{}a policeman\textquotedblright{}.

\begin{grammarlens}[title={What you've built}]
One more word for this chapter.
\end{grammarlens}

\subsection*{Your turn}
\begin{itemize}
\raggedright
  \item \emph{Say it:} \emph{ārakṣakaḥ}
  \item \emph{Say it:} \emph{ārakṣakaḥ}, once more
  \item \emph{Say it:} say \emph{ārakṣakaḥ}
  \item \emph{Recall:} say the Sanskrit for an engineer, then the Sanskrit for a scientist, then say \emph{ārakṣakaḥ} again
\end{itemize}

\begin{tcolorbox}[breakable,colback=teal!4,colframe=teal!35!black,title={Before you move on}]
What does \sk{आरक्षकः} mean? (\textquotedblleft{}A policeman\textquotedblright{}.) Say it once more.
\end{tcolorbox}
\section[\texorpdfstring{\sk{चालकः} (cālakaḥ)}{चालकः (cālakaḥ)}]{\sk{चालकः} (cālakaḥ) — a driver}

\label{lesson:SA-C238-calakah}



\begin{quote}
Before the new one: say the Sanskrit for a scientist, then the Sanskrit for a policeman.
\end{quote}

\subsection*{You'll want to know: \sk{चालकः}}
\textbf{\sk{चालकः}} — \emph{cālakaḥ} — \textquotedblleft{}a driver\textquotedblright{}.

\begin{grammarlens}[title={What you've built}]
One more word for this chapter.
\end{grammarlens}

\subsection*{Your turn}
\begin{itemize}
\raggedright
  \item \emph{Say it:} \emph{cālakaḥ}
  \item \emph{Say it:} \emph{cālakaḥ}, once more
  \item \emph{Say it:} say \emph{cālakaḥ}
  \item \emph{Recall:} say the Sanskrit for a scientist, then the Sanskrit for a policeman, then say \emph{cālakaḥ} again
\end{itemize}

\begin{tcolorbox}[breakable,colback=teal!4,colframe=teal!35!black,title={Before you move on}]
What does \sk{चालकः} mean? (\textquotedblleft{}A driver\textquotedblright{}.) Say it once more.
\end{tcolorbox}
